\section{Related Work}

\paragraph{AI-generated image detection.}
Whole-image AI-generated image detectors usually assume that the entire image carries generator-specific traces. GenImage popularized large-scale evaluation over multiple generators and degradations, but its task is image-level classification rather than localization \citep{zhu2023genimage}. DIRE detects diffusion-generated images through reconstruction discrepancies \citep{wang2023dire}, while CLIP-based variants such as C2P-CLIP improve cross-category generalization \citep{tan2024c2pclip}. These methods are relevant because they are widely deployed as passive AI-image detectors. They are also structurally mismatched to CLAIMFORGE: if only a tiny object region is synthetic, the global image statistics are dominated by real pixels, and an image-level score alone cannot tell an adjudicator where the suspicious evidence lies.

\paragraph{Image manipulation detection and localization.}
Classical image manipulation localization asks for a tamper mask, often together with an image-level authenticity score. TruFor combines localization, an integrity score, and reliability maps for trustworthy forgery analysis \citep{guillaro2022trufor}. HiFi-IFDL studies fine-grained localization across manipulation types and generative edits \citep{guo2023hifi}. IMDL-BenCo unifies models, metrics, and robustness protocols for image manipulation detection and localization \citep{ma2024imdldbenco}. These tools provide the most natural baseline family for CLAIMFORGE because the target output is a forged-region mask. However, many manipulation localizers exploit boundary, noise, compression, or resampling traces. CLAIMFORGE deliberately reduces those shortcuts by preserving the source image and by requiring matched post-processing.

\paragraph{Diffusion inpainting benchmarks.}
GIM and OpenSDI move toward diffusion editing and inpainting-based manipulations at large scale \citep{chen2024gim,wang2025opensdi}. COCO-Inpaint further focuses on detecting and localizing inpainting-based manipulations \citep{yan2025cocoinpaint}. These datasets establish that diffusion-inpainting detection should include both image-level and pixel-level metrics, and they motivate cross-generator evaluation. CLAIMFORGE differs in domain and threat model: it centers on consumer-claim scenes such as restaurants, rooms, bathrooms, and kitchens, and on small object insertions whose semantics directly support a complaint.

\paragraph{Dataset shortcuts and laundering.}
Recent audits show that AI-image detectors can learn construction shortcuts. ``Fake or JPEG?'' demonstrates that real and fake image corpora often differ in storage format, causing detectors to learn format artifacts rather than generation artifacts \citep{grommelt2024fakejpeg}. Calibration work further cautions that fixed thresholds can misrepresent detector separability \citep{yang2026calibration}. Inpainting Exchange gives a direct warning for CLAIMFORGE: detectors that perform well on whole-image generation can degrade toward chance when inpainting removes or launders global traces \citep{nebioglu2026inpaintingexchange}. CLAIMFORGE takes these findings as design constraints: matched post-processing, local-object edits, laundering tests, and AUC/AP reporting are part of the benchmark rather than optional robustness checks.

\paragraph{Segment-anything style forgery localization.}
Recent localization models such as Mesorch and SAFIRE pursue stronger masks and promptable forged-region segmentation \citep{zhu2024mesorch,kwon2024safire}. They are important baselines because they represent current progress on pixel-level forgery localization. CLAIMFORGE is not positioned as a replacement for these models. Instead, it tests whether such models can generalize to claim-relevant, small, semantically plausible insertions under unseen editor, unseen domain, and laundering shifts.
